\section{Discussion}
\label{sec:discussion}

\paragraph{What an uploader can gain.}
On the simulated platforms, the tier choice dominates and is confirmed by every judge. The filter
gain on top of it is small, depends on the platform, and is mostly not confirmed by metrics that were
not optimised. The only filter that survives two of three judges is plain sharpening.
\pending{update with the P1 verdict: whether a learned pre-compensation adds anything for a defined
content class, and with the real-platform results.}

\paragraph{Evaluation lessons.}
Three practices changed our conclusions and, we believe, should be standard for pre-processing work:
decomposing the gain into the tier gain and the filter gain; judging with metrics that the optimiser
never sees; and reporting variability across seeds and clips rather than single runs. Without them,
the composite results alone would have read as a consistent filter gain of about three points over
the tier trick on every natural clip.

\paragraph{Limitations.}
\begin{itemize}
  \item Most results use simulated platforms whose ladders are plausible but invented.
        \pending{state what the real-platform probe confirmed or contradicted.}
  \item Clips are 3\,s long and 720p. Natural content comes from professionally shot test
        sequences, not phone-recorded user content.
  \item Class C has a single clip, so the hypothesis about flat and dark content is untested.
        Confirmation splits contain six or seven clips per class, which limits statistical power.
  \item Retiming 24/25\,fps sources to 30\,fps speeds up their motion and raises TI.
  \item The composite objective's weights were set by hand.
  \item No metric is a substitute for viewers. \pending{viewer study.}
\end{itemize}
